\section{Escalar la montaña nevada sin fin: el camino interior de emprender en IA}

\subsection{De «avanzar hacia una montaña nevada interminable y desconocida» a «Beginning of Infinity»}

Yang Zhilin (杨植麟) es fundador y director ejecutivo de Moonshot AI (月之暗面). Un año después (julio de 2025), en una nueva entrevista, describió lo que sentía recurriendo a la filosofía de \textit{The Beginning of Infinity}:

\begin{importantbox}{Los problemas son inevitables, pero se pueden resolver}
David Deutsch propone en \textit{The Beginning of Infinity} dos tesis centrales:
\begin{enumerate}
    \item Los problemas son inevitables (Problems are inevitable)
    \item Los problemas se pueden resolver (Problems are soluble)
\end{enumerate}
El proceso de investigación y desarrollo de IA encarna precisamente esta filosofía: se resuelven muchos problemas del aprendizaje por refuerzo y surgen problemas nuevos de evaluación y verificación; se resuelve el problema de la eficiencia del preentrenamiento y surge el problema de la generalización de los agentes. Cada vez que se resuelve un problema, la tecnología avanza unos cientos de metros más y, al mismo tiempo, aparecen problemas nuevos.
\end{importantbox}

Yang Zhilin considera que la AGI tal vez no sea un escalón determinado, sino una dirección. La IA actual ya hace mejor que el 99\% de los seres humanos en muchos campos; en cierto sentido puede considerarse una AGI parcial. Pero es difícil decir que en un momento dado se alcanza de pronto la AGI: se parece más a un proceso de ascenso continuo.

\begin{knowledgebox}{Puede que la montaña nevada no tenga fin}
Yang Zhilin expresó una postura optimista: espera que esta montaña nevada nunca tenga fin: «eso es justo lo que significa Beginning of Infinity, puede que sea una montaña infinita». Y en cierta etapa, tal vez ya no sean los seres humanos quienes escalen, sino que se use la IA para escalar; por ejemplo, usar el modelo K2 para participar en el entrenamiento de modelos y en el procesamiento de datos. La IA se está convirtiendo en un «amplificador de la civilización humana».
\end{knowledgebox}

\subsection{Resumen de la sección}

Emprender en IA es un proceso continuo de resolver problemas y descubrir problemas nuevos. La AGI no es un punto final, sino una dirección. La IA está pasando de ser una herramienta a ser un participante en el proceso de investigación y desarrollo.

